\documentclass[12pt]{article}
\usepackage[utf8]{inputenc}
\usepackage[T1]{fontenc}
\usepackage[margin=1in]{geometry}
\usepackage{setspace}
\usepackage{xcolor}
\usepackage[colorlinks=true, urlcolor=blue!60!black, linkcolor=blue!60!black]{hyperref}
\setlength{\parindent}{0em}
\setlength{\parskip}{0.45em}
\renewcommand{\baselinestretch}{1.0}
\begin{document}
\begin{center}
{\Large \textit{South-Central Timber Development, Inc.\ v.\ Wunnicke}, 467 U.S.\ 82 (1984)\par}
\vspace{0.3em}
{\normalsize Case Brief --- One-Page Edition\par}
\end{center}

\textbf{Memory Jogger.} Alaska's in-state processing condition on state timber was not authorized by Congress or shielded as market participation (\href{https://supreme.justia.com/cases/federal/us/467/82/}{Justia}).

\textbf{Facts.} Alaska sold state timber on condition that buyers process it in Alaska before export. South-Central Timber, which exported raw logs to Japan, sued. Federal rules restricted export of unprocessed timber from federal lands.

\textbf{Procedural History.} The district court enjoined the condition; the Ninth Circuit reversed.

\textbf{Votes.} Reversed, 6--2: White for the Court and for a four-Justice plurality; Rehnquist and O'Connor dissented.

\textbf{Holding.} Congressional consent must be ``unmistakably clear''; the plurality held that a market participant may not regulate a downstream market.

\textbf{Reasoning.} A federal-lands policy is not consent for state lands. Alaska participated in timber sales, not processing, and controlled buyers after the sale, unlike \textit{Reeves} and \textit{Alexandria Scrap}; its local-processing rule was virtually \textit{per se} invalid. Rehnquist called the line ``both artificial and unconvincing.''

\textbf{Examples.} \emph{Same-Side:} lithium leases requiring in-state refining. \emph{Opposite-Side:} a state refinery choosing local customers. \emph{Fence-Sitter:} resident logging crews.

\textbf{Critique.} The plurality defined the market by ``intuition'' and drew only four votes, though its line pushes States toward visible subsidies rather than hidden auction discounts.

\textbf{Key Quotation.} ``There is no talismanic significance to the phrase `expressly stated,' however; it merely states one way of meeting the requirement that, for a state regulation to be removed from the reach of the dormant Commerce Clause, congressional intent must be unmistakably clear.'' --- Justice White, 467 U.S.\ at 91--92.

\end{document}
